\documentclass[reqno]{amsart} 
\usepackage{amssymb,latexsym,amsmath,amscd,graphicx,setspace,amsthm,verbatim,enumitem}
\usepackage[margin = 3 cm]{geometry}

\def\upint{\mathchoice%
    {\mkern13mu\overline{\vphantom{\intop}\mkern7mu}\mkern-20mu}%
    {\mkern7mu\overline{\vphantom{\intop}\mkern7mu}\mkern-14mu}%
    {\mkern7mu\overline{\vphantom{\intop}\mkern7mu}\mkern-14mu}%
    {\mkern7mu\overline{\vphantom{\intop}\mkern7mu}\mkern-14mu}%
  \int}
\def\lowint{\mkern3mu\underline{\vphantom{\intop}\mkern7mu}\mkern-10mu\int}


\theoremstyle{plain}
\newtheorem{theorem}{Theorem}[section]
\newtheorem{proposition}{Proposition}
\newtheorem{corollary}{Corollary}
\newtheorem{lemma}{Lemma}
\newtheorem{conjecture}{Conjecture}
\newtheorem{question}{Question}
\newtheorem{problem}{Problem}
      
\theoremstyle{definition}
\newtheorem{definition}{Definition}

\newenvironment{solution1}{\paragraph{\emph{Solution $1$}.}}{\hfill$\square$}
\newenvironment{solution2}{\paragraph{\emph{Solution $2$}.}}{\hfill$\square$}
\newenvironment{solution3}{\paragraph{\emph{Solution $3$}.}}{\hfill$\square$}

\begin{document} 

\title[Homework 8]{Homework 8}

\date{\today} 
\maketitle 

Let $(X,\mathcal{A},\mu)$ be a measure space.  In class, we defined what it means for a property to hold \emph{almost everywhere} on $X$.  It means that there exists a measurable set $A \in \mathcal{A}$ such that $\mu(A) = 0$ and for which the property holds true for all $x \in X \smallsetminus A$.

\begin{problem}
Let $(X,\mathcal{A},\mu)$ be a measure space such that $\mu(X)<\infty$ and let $f:X \rightarrow \mathbb{R}$ be a bounded function.  Assume that $f = 0$ almost everywhere on $X$.  Show that $f$ is Lebesgue integrable on $X$ and that
$$\int_{X}f = 0. $$
\end{problem}
\begin{solution1}

\end{solution1}

\begin{problem}
Let $(X,\mathcal{A},\mu)$ be a measure space such that $\mu(X)<\infty$ and let $g,f:X \rightarrow \mathbb{R}$ be bounded functions.  Assume that $f$ is Lebesgue integrable.  Show that if $f = g$ almost everywhere on $X$, then $g$ is also Lebesgue integrable and
$$\int_{X}g = \int_{X}f. $$
\end{problem}
\begin{solution1}

\end{solution1}

The rest of this homework set is about measurable functions for which we now recall the definition from class.  If $(X,\mathcal{A})$ is a measurable space, then a function $f:X \rightarrow \mathbb{R}$ is called measurable if for all $\lambda \in \mathbb{R}$, one has
$$\{x \in X : f(x) > \lambda\} \in \mathcal{A}. $$

\begin{problem}
Let $(X,\mathcal{A})$ be a measurable space and let $f:X \rightarrow \mathbb{R}$ be a constant function meaning that there exists $c \in \mathbb{R}$ such that $f(x) = c$ for all $x \in X$.  Show that $f$ is a measurable function.
\end{problem}
\begin{solution1}

\end{solution1}


\begin{problem}
Let $(X,\mathcal{A})$ be a measurable space and let $f,g:X \rightarrow \mathbb{R}$ be measurable function.  Show that $f + g$, $f - g$, $f \cdot g$, and $f/g$ provided $g(x) \neq 0$ for all $x \in X$, are all measurable functions.
\end{problem}
\begin{solution1}

\end{solution1}


\begin{problem}
Let $D \in \mathcal{L}(\mathbb{R}^{n})$.  We view $D$ as a measurable space with the $\sigma$-algebra given by the restriction of $\mathcal{L}(\mathbb{R}^{n})$.  Show that a continuous function $f: D \rightarrow \mathbb{R}$ is measurable.
\end{problem}
\begin{solution1}

\end{solution1}

\begin{problem}
Let $D \in \mathcal{L}(\mathbb{R})$.  We view $D$ as a measurable space with the $\sigma$-algebra given by the restriction of $\mathcal{L}(\mathbb{R})$.  Show that if a function $f:D \rightarrow \mathbb{R}$ is increasing, then $f$ is measurable.  (Remember that being increasing means $x \le y$ implies $f(x) \le f(y)$.)
\end{problem}
\begin{solution1}

\end{solution1}




\end{document}









